% 14-cognition.tex: Атлас, Афина и слои, которым никогда не дозволено моделировать человека.
\section{Как Полярис наблюдает и понимает сам себя}
\label{sec:cognition}

Два познающих слоя существуют, два предложены, и одна черта отделяет все их от людей, которым
система служит. Проектное правило записано в репозитории одной фразой: делать читаемой власть, а
не людей.

\figfile{cognition}{Познающие слои. Атлас видит, что делает система; Афина понимает устройство
власти; Метида, пока лишь предложенная, изучала бы поведение самого Поляриса; Мирмекс, предложенный
в настоящем документе, нападал бы на него снаружи. Конституция стесняет их все. Красная черта есть
запрет на человека: ничто над ней не моделирует физическое лицо, и
\code{проверка\_афина\_без\_человека} роняет сборку, если в слое власти появится поле или обход на
уровне человека.}{fig:cognition}

\subsection{Атлас видит}

Атлас есть поверхность расследования для оператора: аналитическая консоль, которая открывается
ограниченным обзором без географии, временным рядом объёма с наложенными отказами, разбивками
<<первые К>> по контексту, ведомству, уровню раскрытия и исходу, двумерной сводной таблицей,
сеткой записей и картой, оставленной на отдельной вкладке для того, в чём карта сильна: для
разбора одного региона или пути одного субъекта. Всякое представление есть ограниченная сводка,
вычисленная на сервере, и никогда не сами события: самое большее пять тысяч сводок кластеров
на область обзора, какую бы мелкую сетку ни запросил клиент, самое большее 240 временных корзин, самое большее пятьдесят категорий и
постраничная выдача по ключу, чтобы глубокая страница оставалась дешёвой. Проверка с нулевым
разглашением учитывается в объёме и в подсчётах по раскрытию, но никогда не получает места на
карте, потому что не несёт токена, по которому её можно было бы приписать. Страницы расследования
по одной сущности построены так, чтобы отказывать в анализе связей между людьми, и всякое чтение
поверхности аудита само записывается. Режим живого моделирования водят в безголовом браузере в НИ,
и утверждается, что консоль действительно получает поток. Начиная с \code{в9.402} замер измеряет
каждую ограниченную сводку на десятой части потока событий и на всём потоке, с разницей в минуты
на одном и том же железе, и оценивает каждую относительно роста числа строк, потому что замер в
одном масштабе говорит лишь, что сводка однажды была быстрой, и не отличает свёртку, идущую вровень
со строками, от свёртки, растущей квадратично; всякая сводка растёт медленнее своих данных, и замер
завершается с ошибкой, если хоть одна их обгонит. \sImpl

\subsection{Афина понимает устройство}

Афина есть слой смысла и происхождения, работающий только на чтение, поверх таблиц власти:
ведомств, алгоритмов, контекстов, аттестаций доверия и политик хранения, в виде десяти
представлений объектов и восьми представлений рёбер, с четырьмя функциями, которые отвечают на
вопросы управления, на какие иначе дерево ответило бы лишь вручную. Почему это ведомство вправе
выдавать под этим алгоритмом? Какая политика доказательства и раскрытия ограничивает этот
контекст? Если этот алгоритм объявить устаревшим, кого это затронет? Какой механизм закрепляет это
конституционное правило? Сама конституция есть данные: С1--С10 и призвание суть строки, а таблица
сопоставляет каждому правилу точный живой механизм, который его закрепляет: триггер, частичный
уникальный индекс, именованное ограничение целостности, функцию проверки или процедуру. Проверка
роняет сборку, если хоть один названный механизм перестал существовать, а тест базы данных
доказывает, что каждый присутствует в работающем каталоге, и это закрывает расхождение между
прозой и кодом, которого не могла закрыть одна лишь решётка ограничений. Подписанный реестр из
раздела~\ref{sec:institutions} есть представление над Афиной, так что карта учреждений, которую
забирает посторонний, есть та же карта, которую рассматривает оператор. \sImpl

Афине безопасно существовать лишь потому, что читаемость людей сделана структурно невозможной, а не
просто нежелательной. Её запирают пять инвариантов, у каждого есть состязательный тест обнаружения:
ни одно представление, функция или таблица Афины не ссылается на таблицу или поле физического
лица, и ни одно поле не есть устойчивый суррогат человека; всякая функция работает только на
чтение, так что умозаключение \emph{граф говорит отозвать, значит отозвано} невозможно; всякая
функция ограничивает свой вывод; всякое ребро власти разрешается в существующую таблицу власти, а
текущие представления исключают отозванные полномочия, так что Афина описывает и никогда не
распоряжается; и всякое правило сопоставлено существующему механизму. Тот же выпуск, что создал
Афину, убрал прежние представления, сводившие данные по людям. Если запрет когда-нибудь помешает
какой-то возможности, собственная оценка репозитория говорит, что это сигнал остановиться, а не
ослабить запрет. \sImpl

\subsection{Чего здесь не должно быть построено никогда}

Никакого объекта человека, домохозяйства или отношения. Никакого дескриптора субъекта,
непрозрачного или какого угодно, устойчивого между контекстами. Никакого обхода от события с
нулевым разглашением к субъекту. Никакого деления населения на сегменты и никаких оценок. Никакого
действия: слой читает и никогда не пишет. Таковы критерии прекращения из того проектного
исследования, которое одобрило стеснённую Афину, и такова граница, которую унаследует всякий
будущий познающий слой.

\subsection{Метида и Фемида, названные осторожно}

Дорожная карта несёт записанное задание на будущую подсистему обучения, которая лишь советует и
изучала бы сам Полярис: обнаружение отклонений в инфраструктуре, прогноз ёмкости с оценкой
неопределённости, разбор регрессий и рисков выпуска, обучение на хаосе, которое предлагает опыты и
никогда их не проводит само, помощь в поиске первопричины по графу зависимостей Афины и сводные
отклонения в работе учреждений вида \emph{это ведомство обычно отзывает 0,8~процента в месяц, а
теперь отзывает 4,9}. Её главную гипотезу надлежит опровергнуть прежде, чем появится хоть строка
кода: машинное обучение вправе выводить свойства системы и никогда не вправе выводить, ранжировать,
оценивать или предсказывать права физических лиц, и переименование человека в субъекта этому
требованию не удовлетворяет, если данные остаются связываемыми. Её проектный закон: только совет.
Наблюдать, моделировать, предсказывать, объяснять, рекомендовать, и никогда не предсказывать, а
затем отзывать, отказывать, выдавать или менять конституцию либо промышленную настройку. Задание
предлагает имя Метида, потому что очевидная альтернатива сталкивается со средствами метрик, и
записывает, что правило, работающее не хуже модели, следует предпочесть модели. Ничего из этого не
построено; по операционному договору из раздела~\ref{sec:research} и не будет построено, пока в этом
не возникнет нужда у внешней стороны. \sFut

Настоящий документ пользуется именем Фемида для роли конституции, стесняющей все четыре слоя, и
собственное задание репозитория употребляет его в том же предложении: Атлас видит, Афина понимает
устройство, Метида изучает поведение системы, Фемида стесняет их всех. Фемида не есть составная
часть. Это призвание и С1--С10 в том виде, как они закреплены в схеме и удерживаются проверками,
которые уже существуют. Имя есть подспорье для чтения, и ничто в коде его не носит.

\begin{layercard}{Карточка слоя: познание системы}
\qa{Что это}{Атлас, ограниченная эксплуатационная консоль; Афина, слой власти и конституции, работающий только на чтение; и два предложения, Метида и Мирмекс.}
\qa{Зачем оно существует}{Оператор обязан видеть, что делает система, и понимать, откуда берутся полномочия; системой такого размера невозможно управлять вручную.}
\qa{Чему доверяет}{Событиям, работающим только на добавление, и таблицам власти, которые читает.}
\qa{Чему отказывает}{Любому объекту, дескриптору, обходу, оценке или действию на уровне человека; неограниченному чтению; сводке, растущей быстрее своих данных; модели там, где справилось бы правило.}
\qa{Что видит}{Сводки событий и устройство власти, ключей, контекстов, доверия и правил.}
\qa{Чего не видит}{Человека. Страница расследования по одной сущности существует вне этих слоёв, заперта ролью и проходит аудит при каждом чтении.}
\qa{Если откажет}{Ссылка на человека где угодно в коде СКЛ Афины роняет \code{проверка\_афина\_без\_человека} и останавливает сборку, что измерено добавлением такой ссылки (прочие проверки Афины охраняют другие свойства и остаются зелёными); неограниченная сводка роняет С8; обучающая система, которая оценивала бы людей, нарушила бы призвание и отвергается прежде, чем её спроектируют.}
\qa{Кем проверяется}{\code{проверка\_афина\_без\_человека}, \code{проверка\_афина\_только\_чтение}, \code{проверка\_функции\_афины\_ограничены}, \code{проверка\_афина\_не\_распоряжается}, \code{проверка\_исполнение\_правил\_афины\_разрешается}, \code{проверка\_консоль\_атласа}, \code{проверка\_замер\_измеряет\_рост}, потолки С8.}
\qa{Сейчас или потом}{Атлас и Афина сейчас. Метида после оценки и лишь если в ней нуждается внешняя сторона; Мирмекс есть предложение настоящего документа.}
\qa{Внутри и снаружи}{Внутри: эксплуатация. Снаружи: петля, которой система исправляет собственное понимание.}
\end{layercard}

\nextlayer{Система, способная наблюдать себя, всё равно должна учиться на том, что наблюдает, а
самый трудный урок, который выучил Полярис, состоит в том, что его собственные инструменты могут
быть зелёными, измеряя не то. Следующий раздел есть этот урок, и настоящей версии есть что к нему
добавить.}
